% -----------------------------------------------------------------------------
% Section 7: Conclusion
% Paper: "Silent Decisions Are Type Errors: Enforcing AI Accountability
%         via Linear Resource Types"
% Venue: IEEE Computer — Special Issue on AI Governance
% Deadline: 13 April 2026
% Authors: Daniel Lau Pereira Soares
% -----------------------------------------------------------------------------

\section{Conclusion}
\label{sec:conclusion}

This paper began from a simple observation: regulatory frameworks
such as LGPD Art.~20 and EU AI Act Arts.~14 and~86 establish review,
information, explanation, and oversight obligations, but do not by
themselves prevent software from emitting a decision without supporting
evidence.
We called this failure mode a \emph{silent decision}, and argued
that its root cause is not a failure of intent but a failure of
type---a structural property of architectures that treat
accountability as a runtime audit obligation rather than a
compile-time invariant.

Drawing on Girard's Linear Logic~\cite{girard1987linear} and Rust's
ownership model~\cite{jung2017rustbelt}, we introduced the
Bound-To-Verdict (BTV) framework and its core type law:
$(E \otimes C) \multimap V$.
Under this invariant, a \texttt{Verdict} cannot be constructed
without atomically consuming an \texttt{EvidenceToken} and a
\texttt{ComplianceToken}.
The structural protections---private fields, \texttt{pub(crate)}
consumption, and an affine ownership API reinforced by
\texttt{\#[must\_use]} lint policy---enclose the authorized constructor.
Within that perimeter, an unevidenced BTV verdict construction is rejected
at compile time.

We proved the \emph{Constitutional Enclosure Theorem}: under the
assumptions of Section~4.3, an external Safe Rust caller cannot obtain
a BTV \texttt{Verdict} through the public API without the required
tokens.  The construction paths are machine-checked
(\texttt{cargo test}: seventeen clauses, seconds), with a scoped
Fail-Secure Corollary: the protected path cannot issue a verdict
without evidence.

Empirical benchmarks on the reference implementation bound what the
type-level enforcement itself costs: it adds no runtime branch.  The
measured cost is that of the cryptographic primitives and encoding
(BLAKE3 + two HMAC-SHA256 computations, 5.36--7.65\,\textmu s for the
full public construction path across 64\,B--4\,KiB payloads) and of durable
persistence.  On the reduced-footprint reference collection, durable SQLite
reached a 261\,\textmu s median at one thread and a 28\,ms estimated p99 at
four threads.  These are container-storage measurements from one current
x86-64 host, not production-scale or cross-architecture performance claims.

The implications extend beyond compliance engineering.
If accountability can be expressed as a type invariant, then
the correctness of governance becomes a property that compilers
can check, proof assistants can verify, and regulators can audit
through \texttt{cargo test} rather than through manual code
review.
BTV therefore prevents a class of unevidenced verdict constructions
inside a defined type perimeter.  It does not prove that every
consequential action in a polyglot system is routed through that perimeter.
